\documentclass[../../main.text]{subfiles}
%Deutsche Sprache
\usepackage[ngerman]{babel}
\usepackage[T1]{fontenc}
\usepackage{lmodern}
\usepackage[thmmarks,amsmath,hyperref,noconfig]{ntheorem}

%Zeilenabstand
\usepackage[onehalfspacing]{setspace}

\begin{document}
\begin{doublespace}
%%%%%%%%%%%%%%%%%%%%%%%%%%%%%%%%%%%%%%%%%%%%%%%%
\subsection{Large-Language Modelle}
Large Language Modelle beschreibt ein Modell, welches aufgrund von
mathematischen Berechnungen, Text erzeugen kann. [Quelle]
Ein solches Modell hat durch Training erlernt auf Grundlage von
Statistik, Wortfolgen zu erstellen. Dieses Wissen und diese
Möglichkeit zur Erzeugung wird genutzt um mit dem User zu interagieren.
Einsatzgebiete sind unter anderem die Erzeugung von Code, das Erstellen von
Schriftstücken aller Art so auch Gedichte oder Geschichten oder auch
zum Beweis von mathematischen Problemen (so soll angeblich ein 
\href{https://www.spiegel.de/wissenschaft/kuenstliche-intelligenz-ki-von-openai-soll-millenniumsraetsel-in-88-stunden-geloest-haben-a-ae808aee-1a76-4ad8-9a90-d870f4d64892}{Millenium-Problem} von
OpenAI gelöst worden sein).

%%%%%%%%%%%%%%%%%%%%%%%%%%%%%%%%%%%%%%%%%%%%%%%%
\subsection{Systemkonfiguration}
%%%%%%%%%%%%%%%%%%%%%%%%%%%%%%%%%%%%%%%%%%%%%%%%
\subsection{Halluzinationen}
%%%%%%%%%%%%%%%%%%%%%%%%%%%%%%%%%%%%%%%%%%%%%%%%
\subsection{Vertrauen in Aussagen}
%%%%%%%%%%%%%%%%%%%%%%%%%%%%%%%%%%%%%%%%%%%%%%%%
\end{doublespace}
\end{document}